% Options for packages loaded elsewhere
\PassOptionsToPackage{unicode}{hyperref}
\PassOptionsToPackage{hyphens}{url}
\documentclass[
]{article}
\usepackage{xcolor}
\usepackage[margin=1in]{geometry}
\usepackage{amsmath,amssymb}
\setcounter{secnumdepth}{-\maxdimen} % remove section numbering
\usepackage{iftex}
\ifPDFTeX
  \usepackage[T1]{fontenc}
  \usepackage[utf8]{inputenc}
  \usepackage{textcomp} % provide euro and other symbols
\else % if luatex or xetex
  \usepackage{unicode-math} % this also loads fontspec
  \defaultfontfeatures{Scale=MatchLowercase}
  \defaultfontfeatures[\rmfamily]{Ligatures=TeX,Scale=1}
\fi
\usepackage{lmodern}
\ifPDFTeX\else
  % xetex/luatex font selection
\fi
% Use upquote if available, for straight quotes in verbatim environments
\IfFileExists{upquote.sty}{\usepackage{upquote}}{}
\IfFileExists{microtype.sty}{% use microtype if available
  \usepackage[]{microtype}
  \UseMicrotypeSet[protrusion]{basicmath} % disable protrusion for tt fonts
}{}
\makeatletter
\@ifundefined{KOMAClassName}{% if non-KOMA class
  \IfFileExists{parskip.sty}{%
    \usepackage{parskip}
  }{% else
    \setlength{\parindent}{0pt}
    \setlength{\parskip}{6pt plus 2pt minus 1pt}}
}{% if KOMA class
  \KOMAoptions{parskip=half}}
\makeatother
\usepackage{longtable,booktabs,array}
\usepackage{calc} % for calculating minipage widths
% Correct order of tables after \paragraph or \subparagraph
\usepackage{etoolbox}
\makeatletter
\patchcmd\longtable{\par}{\if@noskipsec\mbox{}\fi\par}{}{}
\makeatother
% Allow footnotes in longtable head/foot
\IfFileExists{footnotehyper.sty}{\usepackage{footnotehyper}}{\usepackage{footnote}}
\makesavenoteenv{longtable}
\usepackage{graphicx}
\makeatletter
\newsavebox\pandoc@box
\newcommand*\pandocbounded[1]{% scales image to fit in text height/width
  \sbox\pandoc@box{#1}%
  \Gscale@div\@tempa{\textheight}{\dimexpr\ht\pandoc@box+\dp\pandoc@box\relax}%
  \Gscale@div\@tempb{\linewidth}{\wd\pandoc@box}%
  \ifdim\@tempb\p@<\@tempa\p@\let\@tempa\@tempb\fi% select the smaller of both
  \ifdim\@tempa\p@<\p@\scalebox{\@tempa}{\usebox\pandoc@box}%
  \else\usebox{\pandoc@box}%
  \fi%
}
% Set default figure placement to htbp
\def\fps@figure{htbp}
\makeatother
\setlength{\emergencystretch}{3em} % prevent overfull lines
\providecommand{\tightlist}{%
  \setlength{\itemsep}{0pt}\setlength{\parskip}{0pt}}
\usepackage{bookmark}
\IfFileExists{xurl.sty}{\usepackage{xurl}}{} % add URL line breaks if available
\urlstyle{same}
\hypersetup{
  pdftitle={JIJOMAO},
  pdfauthor={Carlos Fernando Ramírez Dávila},
  hidelinks,
  pdfcreator={LaTeX via pandoc}}

\title{JIJOMAO}
\author{Carlos Fernando Ramírez Dávila}
\date{2026-04-06}

\begin{document}
\maketitle

\begin{enumerate}
\def\labelenumi{\Roman{enumi}.}
\setcounter{enumi}{1}
\tightlist
\item
  Algoritmos de Identificación: Matemática Completa
\end{enumerate}

II.1 --- Modelo base: VAR(p) reducido

El sistema estimado es un VAR de forma reducida de orden \(p=3\) con
bloque exógeno (Cushman \& Zha, 1997):

\[\mathbf{y}t = \mathbf{c} + \sum{l=1}^{p} \mathbf{B}l,\mathbf{y}{t-l} + \mathbf{D},\mathbf{z}_t + \mathbf{u}_t, \qquad \mathbf{u}_t \sim
  \text{iid}\bigl(\mathbf{0},,\boldsymbol{\Sigma}_u\bigr)\]

donde \(\mathbf{y}_t \in \mathbb{R}^6\) son las variables endógenas y
\(\mathbf{z}_t \in \mathbb{R}^3\) son las exógenas (bloque EE.UU.). La
estimación es OLS\\
ecuación por ecuación (B = (X'X)(X'Y) --- línea 99 del MATLAB).

El SVAR estructural asociado es:

\[\mathbf{A}_0,\mathbf{y}t = \tilde{\mathbf{c}} + \sum{l=1}^{p}\mathbf{A}l,\mathbf{y}{t-l} + \boldsymbol{\varepsilon}_t, \qquad \boldsymbol{\varepsilon}_t \sim   
  \text{iid}(\mathbf{0},,\mathbf{I}_n)\]

La relación entre formas es
\(\mathbf{u}_t = \mathbf{A}_0^{-1}\boldsymbol{\varepsilon}_t\), por lo
que:

\[\boldsymbol{\Sigma}_u = \mathbf{A}_0^{-1}\bigl(\mathbf{A}_0^{-1}\bigr)' \equiv \mathbf{B},\mathbf{B}', \qquad \mathbf{B} \equiv \mathbf{A}_0^{-1}\]

La función de impulso-respuesta al shock \(j\) en el horizonte \(h\) es:

\[\text{IRF}_{ij}(h) = \mathbf{e}_i',\boldsymbol{\Psi}_h,\mathbf{b}_j\]

donde \(\boldsymbol{\Psi}_h = \mathbf{J},\mathbf{F}^h,\mathbf{J}'\) es
la IRF de forma reducida, \(\mathbf{F}\) es la matriz compañera del VAR,
\(\mathbf{J}\) selecciona las primeras \(n\) filas, y \(\mathbf{b}_j\)
es la \(j\)-ésima columna de \(\mathbf{B}\) (vector de impacto).

\begin{longtable}[]{@{}
  >{\raggedright\arraybackslash}p{(\linewidth - 0\tabcolsep) * \real{0.0556}}@{}}
\toprule\noalign{}
\begin{minipage}[b]{\linewidth}\raggedright
II.2 --- Estrategia 1: Restricciones de Exclusión (Cholesky/Sims)
\end{minipage} \\
\midrule\noalign{}
\endhead
\bottomrule\noalign{}
\endlastfoot
II.3 --- Estrategia 2: Restricciones de Signo Estándar (Uhlig, 2005) \\
Identificación: El conjunto de matrices de impacto válidas es \\
\(\mathcal{B} = \bigl{\mathbf{B} = \mathbf{P},\mathbf{Q} : \mathbf{Q}\in O(n),; \text{sgn}(\mathbf{b}_j) = \mathbf{s}\bigr}\) \\
donde \(\mathbf{s} \in {-1,0,+1}^n\) son las restricciones de signo
(Tabla 2 de C\&E) y \(O(n)\) es el grupo ortogonal. La factorización
\(\mathbf{P}\mathbf{Q}(\mathbf{P}\mathbf{Q})' = \mathbf{P}\mathbf{Q}\mathbf{Q}'\mathbf{P}' = \boldsymbol{\Sigma}_u\)
garantiza que todas las rotaciones son igualmente admisibles en términos
de la forma reducida. \\
Distribución uniforme sobre \(O(n)\) --- medida de Haar: Se genera
\(\mathbf{Q}\) muestreando uniformemente sobre \(O(n)\) via
descomposición QR de una matriz normal aleatoria: \\
\(\mathbf{Z} \sim \mathcal{N}(\mathbf{0},\mathbf{I}_{n\times n}) \quad \Rightarrow \quad \mathbf{Z} = \mathbf{Q},\mathbf{R} \quad \Rightarrow \quad \mathbf{Q}
\leftarrow \mathbf{Q}\cdot\text{diag}\bigl(\text{sgn}(\text{diag}(\mathbf{R}))\bigr)\) \\
El ajuste de signo en la última línea asegura distribución uniforme
sobre \(O(n)\) (no solo sobre \(SO(n)\)) --- ver función
haar\_random\_orthogonal en SVAR\_SignRestrictions.m líneas 192--199. \\
\% SVAR\_SignRestrictions.m --- líneas 282-302 Q =
haar\_random\_orthogonal(n); \% {[}Q,R{]}=qr(randn(n));
Q=Q\emph{diag(sign(diag(R))) A0 = P\_b } Q; resp = A0(:, j\_mp); \%
impacto contemporáneo al choque j \\
valid = true; for vv = 1:n if signs\_mp(vv) \textasciitilde= 0 if
signs\_mp(vv) * resp(vv) \textless= 0 valid = false; break; end end end
if \textasciitilde valid, continue; end \\
Las restricciones impuestas para el choque monetario son (signs\_mp =
{[}-1;-1;0;-1;+1;-1{]}): \\
┌──────────────────────────────┬───────┬──────────────────────────────────────────┐
│ Variable │ Signo │ Razón económica │
├──────────────────────────────┼───────┼──────────────────────────────────────────┤
│ Output (\(y\)) │ \(-\) │ Choque contractivo reduce actividad │
├──────────────────────────────┼───────┼──────────────────────────────────────────┤
│ Inflación subyacente (\(\pi\)) │ \(-\) │ Price puzzle eliminado por
construcción │
├──────────────────────────────┼───────┼──────────────────────────────────────────┤
│ PPI (\(\pi^\text{ppi}\)) │ libre │ Incertidumbre en transmisión │
├──────────────────────────────┼───────┼──────────────────────────────────────────┤
│ Dep. TCR │ \(-\) │ Apreciación real tras alza de tasa (UIP) │
├──────────────────────────────┼───────┼──────────────────────────────────────────┤
│ Tasa de interés (\(i\)) │ \(+\) │ Definición del choque monetario │
├──────────────────────────────┼───────┼──────────────────────────────────────────┤
│ Crecimiento M2 │ \(-\) │ Política restrictiva contrae liquidez │
└──────────────────────────────┴───────┴──────────────────────────────────────────┘ \\
\end{longtable}

II.4 --- Estrategia 3: Restricciones de Signo Aumentadas --- Algorithm 1
de C\&E

Motivación: Las restricciones de signo estándar admiten modelos donde la
inflación responde marginalmente al choque monetario (casi cero), lo
cual es inconsistente con la elasticidad precio-interés estimada en la
economía mexicana. C\&E (2015) proponen una cota de elasticidad como
filtro adicional.

Ecuación 11 --- Bound de elasticidad: A partir de la ecuación de curva
de Phillips forward-looking restringida:

\[\hat{\beta}^- = \hat{\beta}i + \hat{\beta}{ii}, \qquad \hat{\beta}_{lower} = \hat{\beta}^- - 1.645\cdot\widehat{\text{SE}}(\hat{\beta}^-)\]

donde \(\hat{\beta}i\) y \(\hat{\beta}{ii}\) son coeficientes de la
regresión de inflación sobre tasa de interés rezagada e interacción. En
nuestra estimación\\
(Newey-West, 3 lags):

\[\hat{\beta}^- = -0.400, \quad \widehat{\text{SE}} = 0.192 \quad \Rightarrow \quad \hat{\beta}_{lower} = -0.715\]

Restricción aumentada: Para cada candidato
\(\mathbf{A}_0 = \mathbf{P}_b\mathbf{Q}\), además de los signos, se
exige:

\[\hat{\beta}_{lower}\cdot\Delta i ;\leq; \Delta\pi ;\leq; 0\]

donde \(\Delta i = \mathbf{r}{j\text{mp}}\) y
\(\Delta\pi = \mathbf{r}_2\) son las respuestas de impacto de la tasa y
la inflación en la columna del choque monetario. En código:

\% SVAR\_SignRestrictions.m --- líneas 307-312 if strcmp(type,
`augmented') dpi\_imp = resp(2); \% Δπ (inflación core, posición 2)
di\_imp = resp(j\_mp); \% Δi (tasa de interés, posición 5) if
beta\_lower\_val * di\_imp \textgreater{} dpi\_imp continue; \%
rechazar: respuesta de inflación demasiado débil end end

Algorithm 1 (C\&E) --- Descripción completa:

Inputs: Y (T×6), Z (T×3), p=3, B\_boot=2000, K=100, K\_max=5000, h=16
signs\_mp, β\_lower, j\_mp=5

\begin{enumerate}
\def\labelenumi{\arabic{enumi}.}
\item
  Estimar VAR reducido → (B̂, Σ̂\_u, ê\_t) P₀ = chol(Σ̂\_u) ← Cholesky base
\item
  Para b = 1,\ldots,2000: 2a. Resamplear residuos (Runkle 1987): idx =
  randi(T\_eff, T\_eff, 1) ẽ\_b = ê{[}idx, :{]} 2b. Pseudo-datos
  recursivos: Ỹ\_b = runkle\_bootstrap\_data(B̂, Y, Z, p, ẽ\_b) 2c.
  Re-estimar VAR en Ỹ\_b → (B̂\_b, Σ̂\_b) P\_b = chol(Σ̂\_b) ← nueva
  Cholesky del draw

  2d. Para k = 1,\ldots,5000: Q \textasciitilde{} Haar(6×6) ←
  qr(randn(6)): uniforme en O(6) A₀ = P\_b * Q r = A₀{[}:, j\_mp{]} ←
  columna del choque monetario

\begin{verbatim}
Verificar signos:
  ∀ v: signs_mp[v] ≠ 0  →  signs_mp[v] · r[v] > 0  (o rechazar)

(Solo aumentado) Verificar cota:
  β_lower · r[j_mp] ≤ r[2]  (o rechazar)

Si pasa: IRF_full = compute_irf(B̂_b, n, p, r, h)
         Acumular en all_irfs
Si k_acc ≥ K: salir del loop de candidatos
\end{verbatim}
\item
  Sobre all\_irfs (n\_total × n × h+1): IRF\_mediana = median(all\_irfs,
  dim=1) IRF\_lo = prctile(all\_irfs, 16, dim=1) IRF\_hi =
  prctile(all\_irfs, 84, dim=1)
\end{enumerate}

El número total de modelos aceptados depende de qué tan restrictivas son
las condiciones: 198,654 para restricciones estándar y 66,874 para
aumentadas (ratio\\
\textasciitilde3:1), lo que refleja que la restricción de elasticidad
rechaza aprox. el 66\% de los candidatos que pasan los signos.

\begin{longtable}[]{@{}
  >{\raggedright\arraybackslash}p{(\linewidth - 0\tabcolsep) * \real{0.0556}}@{}}
\toprule\noalign{}
\endhead
\bottomrule\noalign{}
\endlastfoot
III. Principales Diferencias: Replicación vs Paper \\
┌─────────────────────────┬──────────────────────────┬───────────────────────────┬──────────────────────┬─────────────────────┐
│ Dimensión │ Paper C\&E │ Replicación │ Magnitud │ Estado │
├─────────────────────────┼──────────────────────────┼───────────────────────────┼──────────────────────┼─────────────────────┤
│ \(T\) │ 159 │ 159 │ Exacto │ ✅ Resuelto │
├─────────────────────────┼──────────────────────────┼───────────────────────────┼──────────────────────┼─────────────────────┤
│ Rezagos (\(p\)) │ 3 │ 3 │ Exacto │ ✅ │
├─────────────────────────┼──────────────────────────┼───────────────────────────┼──────────────────────┼─────────────────────┤
│ Estabilidad VAR │ ✓ │ ✓ │ Exacto │ ✅ │
├─────────────────────────┼──────────────────────────┼───────────────────────────┼──────────────────────┼─────────────────────┤
│ \(\hat{\beta}_{lower}\) │ \(-0.700\) │ \(-0.715\) │ \(\Delta=0.015\)
(\textasciitilde2\%) │ ≈ Concordante │
├─────────────────────────┼──────────────────────────┼───────────────────────────┼──────────────────────┼─────────────────────┤
│ TCR socios │ 21 │ 49 (Banxico SR15997) │ Diferente │ ⚠ Fuente de datos
│
├─────────────────────────┼──────────────────────────┼───────────────────────────┼──────────────────────┼─────────────────────┤
│ Ajuste estacional │ TRAMO/X-12 aditivo │ X-13ARIMA-SEATS │ Metodología
distinta │ ⚠ Impacto menor │
├─────────────────────────┼──────────────────────────┼───────────────────────────┼──────────────────────┼─────────────────────┤
│ Modelos aceptados (std) │ N/D │ 198,654 │ Sin referencia │ --- │
├─────────────────────────┼──────────────────────────┼───────────────────────────┼──────────────────────┼─────────────────────┤
│ Producto doméstico │ Brecha precalculada IGAE │ Índice IGAE nivel
(INEGI) │ Diferente fuente │ ⚠ Resuelto │ │ │ │ │ │ (índice) │
├─────────────────────────┼──────────────────────────┼───────────────────────────┼──────────────────────┼─────────────────────┤
│ IRF signos cualitativos │ Figura 7/8/9 │ Coinciden │ Concordante │ ✅
│
├─────────────────────────┼──────────────────────────┼───────────────────────────┼──────────────────────┼─────────────────────┤
│ Price puzzle │ Eliminado │ Eliminado │ Concordante │ ✅ │
└─────────────────────────┴──────────────────────────┴───────────────────────────┴──────────────────────┴─────────────────────┘ \\
Posibles causas de las diferencias cuantitativas en
\(\hat{\beta}^-\): \\
1. TRAMO vs X-13: diferencias menores en las series de inflación
desestacionalizada afectan la regresión de la Ecuación 11 2. TCR
multilateral vs 49 socios: la composición del índice cambia la dinámica
de dep\_tcr, afectando los coeficientes
\(\hat{\beta}i, \hat{\beta}{ii}\) 3. Brecha del producto: C\&E usa la
brecha IGAE precalculada por Banxico; nosotros extraemos el gap via HP
del índice nivel, lo que produce una serie con tendencia ligeramente
distinta \\
\end{longtable}

\begin{enumerate}
\def\labelenumi{\Roman{enumi}.}
\setcounter{enumi}{3}
\tightlist
\item
  Problema del Signo del TCR bajo Restricciones de Exclusión
\end{enumerate}

El problema original

Al inicio de la replicación, las IRFs de Cholesky mostraban que ante un
choque monetario contractivo (subida de tasa de interés), el TCR se
depreciaba en horizontes medios. Esto es económicamente incorrecto: por
la Paridad de Tasas de Interés (UIP):

\[i_t - i_t^* = \mathbb{E}t[e{t+1}] - e_t \quad \Rightarrow \quad \uparrow i \Rightarrow \text{apreciación del peso} \Rightarrow \text{dep_tcr} < 0\]

Causa raíz: definición de TCR en C\&E vs SR15997

C\&E (2015/2018) definen su TCR multilateral siguiendo la convención
estándar de Banxico en la literatura de los 2000:

\[\text{TCR}_{t}^{\text{C&E}} = \frac{e_t \cdot P_t^*}{P_t}\]

donde un aumento = depreciación real del peso. Bajo esta definición,
dep\_tcr \textgreater{} 0 significa que el peso se deprecia en términos
reales.

La serie SR15997 de SIE Banxico (``Índice de tipo de cambio real
multilateral'') usa la misma orientación: aumento = depreciación. Esta
es la convención del código:

\% SVAR\_SignRestrictions.m --- línea 18 (comentario de cabecera) \% 4:
dep\_tcr --- Crecimiento TCR (+ = depreciación)

\% SVAR\_SignRestrictions.m --- línea 127 (restricciones de signo)
signs\_mp = {[}-1; -1; 0; -1; 1; -1{]}; \% \^{}--- dep\_tcr debe ser
NEGATIVO (apreciación) tras choque contractivo

La causa del signo erróneo

En versiones previas de la replicación, el código de procesamiento
calculaba:

\# VERSIÓN INCORRECTA (invertida): dep\_tcr = (log(lag(tcr)) - log(tcr))
* 1200 \# ← signo invertido

en lugar de:

\# Code/0Proc\_data.R --- línea 185 (VERSIÓN CORRECTA): dep\_tcr =
(log(tcr) - log(lag(tcr))) * 1200 \# + = depreciación

Con la versión invertida, un aumento en SR15997 (depreciación) producía
un dep\_tcr negativo, invirtiendo la interpretación. Cuando la IRF de
Cholesky mostraba\\
apreciación en horizontes \(h > 0\), la serie registraba eso como un
valor positivo (porque la tasa de cambio estaba invertida), resultando
en el signo contrario al esperado en las gráficas.

Cómo está resuelto

En el flujo actual:

\begin{enumerate}
\def\labelenumi{\arabic{enumi}.}
\tightlist
\item
  0Proc\_data.R calcula dep\_tcr = (log(tcr) - log(lag(tcr))) * 1200 →
  convención correcta (positivo = depreciación)
\item
  SVAR\_SignRestrictions.m impone signs\_mp(4) = -1 → el choque
  contractivo debe producir apreciación (dep\_tcr negativo) ✅
\item
  Las IRFs de Cholesky (Fig. 7) muestran dep\_tcr que no reacciona en
  \(h=0\) (por construcción, la variable 4 está antes del choque en la
  posición 5) y que muestra apreciación en horizontes medios,
  consistente con la UIP y con la Figura 7 del paper original.
\end{enumerate}

▎ Nota sobre los 49 vs 21 socios comerciales: aunque este cambio no
afecta el signo, sí cambia la magnitud de las IRFs del TCR. SR15997
pondera por comercio total (49 socios); el índice de C\&E pondera los 21
principales socios de 2001. En ambos casos la apreciación post-choque
monetario es cualitativa y cuantitativamente\\
similar.

\end{document}
